\scanpage{25}
près, ce sont les composantes connexes des
foncteurs de la forme $H_J$.

\uline{Foncteurs spéciaux} $\mathcal E\mathbin{\text{---}}(\mathrm{Ens}_f)$ : ce sont les
foncteurs qui sont \inserted{isom. à des} sommes disjointes finies
de foncteurs de la forme $h_I$ --- ou encore
qui sont des \unsure{sommandes|sommants} directs
(pour $I$ de cardinal assez grand) des
foncteurs de la forme $H_I$.

\uline{Stabilités} Toute somme finie, \inserted{et} produit
fini \inserted{(plus gén. les $\varprojlim$ finies)} de foncteurs spéciaux est spécial.
Tout foncteur spécial $\mathcal E\mathbin{\text{---}}(\mathrm{Ens}_f)$ prend
ses valeurs dans la sous-catégorie
$\mathcal E$ de $(\mathrm{Ens}_f)$, i.e. \deleted{\unsure{donne}}\inserted{définit un} $\in\uline{\operatorname{Hom}}(\mathcal E,\mathcal E)$
(et on identifie par abus de langage
le foncteur $\mathcal E\longrightarrow(\mathrm{Ens}_f)$ avec $\mathcal E\mathbin{\text{---}}\mathcal E$)
Avec cette convention, le composé de
deux foncteurs spéciaux est spécial.
\mspage{24}
